\documentclass[../DoAn.tex]{subfiles}
\begin{document}

Chương này trình bày phần lõi của đồ án: xây dựng các hợp đồng thông minh. Mở đầu là kiến trúc tổng thể, sơ đồ gói và sơ đồ phụ thuộc giữa các hợp đồng, sau đó đi vào từng hợp đồng theo đúng bốn cơ chế đã nêu ở Chương~\ref{chapter:Related_works}. Phần backend và frontend đóng vai trò hỗ trợ và được trình bày ngắn gọn ở cuối chương.

\section{Kiến trúc tổng thể và luồng giao dịch}
\label{sec:arch}

Hệ thống NeutronX gồm bốn tầng. Tầng \textbf{hợp đồng thông minh} chạy trên chuỗi, chứa toàn bộ logic kiểm tra và thực hiện lệnh; đây là phần quan trọng nhất. Tầng \textbf{backend} chạy ngoài chuỗi, đảm nhận việc đồng ký lệnh, lưu trữ và theo dõi trạng thái. Tầng \textbf{filler} là các chương trình thực hiện khớp lệnh. Cuối cùng là \textbf{giao diện người dùng} để tạo và theo dõi lệnh. Hình~\ref{fig:arch} mô tả cách bốn tầng này nối với nhau.

\begin{figure}[H]
    \centering
    \noindent\makebox[\linewidth][c]{\includegraphics[width=1.3\linewidth]{Hinhve/arch.png}}
    \caption{Kiến trúc tổng thể của hệ thống NeutronX gồm bốn tầng.}
    \label{fig:arch}
\end{figure}

Tầng hợp đồng gồm hai nhóm chính: nhóm \textbf{cùng chuỗi} (PartialFillReactor, FillAuction, FallbackExecutor) xử lý lệnh trên một chuỗi duy nhất; nhóm \textbf{xuyên chuỗi} (EscrowSrcFactory, EscrowDstFactory và các hợp đồng ký quỹ per-slot) xử lý lệnh giữa hai chuỗi khác nhau. Cả hai nhóm dùng chung một tập thư viện nội bộ và phụ thuộc vào Permit2 để chuyển tài sản. Hình~\ref{fig:packages} phân nhóm các hợp đồng và thư viện, Hình~\ref{fig:deps} thể hiện quan hệ gọi giữa chúng.

\begin{figure}[H]
    \centering
    \includegraphics[width=0.90\textwidth]{Hinhve/contract_packages.png}
    \caption{Sơ đồ gói các hợp đồng thông minh của NeutronX.}
    \label{fig:packages}
\end{figure}

\begin{figure}[H]
    \centering
    \includegraphics[width=0.95\textwidth]{Hinhve/contract_deps.png}
    \caption{Sơ đồ phụ thuộc giữa các hợp đồng.}
    \label{fig:deps}
\end{figure}

Vòng đời của một lệnh trong cùng một chuỗi diễn ra qua các bước sau:

\begin{enumerate}
    \item Người dùng tạo lệnh trên giao diện và ký lệnh bằng ví theo chuẩn EIP-712. Việc ký này diễn ra ngoài chuỗi nên không tốn phí gas.
    \item Backend nhận lệnh, đồng ký (cosign) để xác nhận các tham số giá, rồi lưu vào cơ sở dữ liệu và công bố cho các filler.
    \item Một filler thấy lệnh có lợi thì đăng ký nhận một phần khối lượng, kèm theo một khoản cọc, thông qua hợp đồng \texttt{FillAuction}.
    \item Filler gọi hàm khớp lệnh trên \texttt{PartialFillReactor}. Hợp đồng dùng Permit2 để chuyển tài sản vào của người dùng cho filler, đồng thời filler trả tài sản ra cho người dùng.
    \item Sau mỗi lần khớp, \texttt{FillAuction} hoàn cọc cho filler theo mức họ đã giao đúng cam kết. Quá trình lặp lại cho tới khi lệnh được khớp hết.
    \item Nếu gần tới hạn mà lệnh vẫn còn dư, \texttt{FallbackExecutor} sẽ khớp nốt phần còn lại.
\end{enumerate}

Hình~\ref{fig:sequence} biểu diễn luồng này dưới dạng sơ đồ trình tự.

\begin{figure}[H]
    \centering
    \noindent\makebox[\linewidth][c]{\includegraphics[width=1.3\linewidth]{Hinhve/sequence_samechain.png}}
    \caption{Sơ đồ trình tự vòng đời một lệnh được khớp từng phần trong cùng một chuỗi.}
    \label{fig:sequence}
\end{figure}

\section{Mô hình dữ liệu lệnh}
\label{sec:order-model}

Mọi thông tin của một lệnh được gói trong một cấu trúc gọi là \texttt{PartialFillOrder}, gồm 11 trường như ở Bảng~\ref{table:order_fields}. Toàn bộ 11 trường này đều nằm trong phần dữ liệu được ký, nên không bên nào có thể tự ý thay đổi mà không bị phát hiện.

\begin{table}[H]
    \centering
    \caption{Các trường của một lệnh \texttt{PartialFillOrder}.}
    \label{table:order_fields}
    \begin{tabular}{|l|l|}
        \hline
        \textbf{Trường} & \textbf{Ý nghĩa} \\
        \hline
        \texttt{swapper}        & Địa chỉ người tạo lệnh. \\
        \hline
        \texttt{inputToken}     & Token người dùng bán ra. \\
        \hline
        \texttt{inputAmount}    & Tổng khối lượng token bán. \\
        \hline
        \texttt{outputToken}    & Token người dùng muốn nhận. \\
        \hline
        \texttt{minOutputAmount}& Mức nhận tối thiểu chấp nhận được (chống trượt giá). \\
        \hline
        \texttt{deadline}       & Số block hạn chót của lệnh. \\
        \hline
        \texttt{nonce}          & Số chống lặp lại lệnh. \\
        \hline
        \texttt{minFillBps}     & Tỉ lệ tối thiểu cho mỗi lần khớp (tránh khớp quá vụn). \\
        \hline
        \texttt{startPrice}     & Giá khởi điểm của đường cong đấu giá Hà Lan. \\
        \hline
        \texttt{decayPerBlock}  & Mức giảm giá sau mỗi block. \\
        \hline
        \texttt{feeTier}        & Mức phí pool dùng để tham chiếu giá. \\
        \hline
    \end{tabular}
\end{table}

Ba trường cuối (\texttt{startPrice}, \texttt{decayPerBlock}, \texttt{feeTier}) mô tả đường giảm giá. Ba trường này nằm trong phần dữ liệu được ký, do đó đường giảm giá được cosigner xác thực và filler không thể tự đặt tham số giá có lợi cho bản thân.

\section{Hợp đồng PartialFillReactor}
\label{subsec:reactor}

\texttt{PartialFillReactor} là hợp đồng đảm nhận việc khớp lệnh từng phần. Bảng~\ref{table:pfr_methods} liệt kê các hàm chính.

\begin{table}[H]
    \centering
    \caption{Các hàm chính của \texttt{PartialFillReactor}.}
    \label{table:pfr_methods}
    \begin{tabular}{|p{4cm}|p{4.5cm}|p{5.5cm}|}
        \hline
        \textbf{Hàm} & \textbf{Tham số / Kiểu trả về} & \textbf{Mục đích} \\
        \hline
        \texttt{register} &
            \texttt{order, sig, cosigSig,} \newline \texttt{fillBps: uint16} &
            Filler cam kết khớp \texttt{fillBps}\% lệnh, đặt cọc qua FillAuction. \\
        \hline
        \texttt{execute}\newline\texttt{PartialChunk} &
            \texttt{order, sig, cosigSig,} \newline \texttt{fillAmount: uint256} &
            Một lần khớp: Permit2 kéo input từ người dùng, filler trả output. \\
        \hline
        \texttt{cancel} &
            \texttt{order} &
            Người dùng hủy lệnh chưa khớp hết. \\
        \hline
        \texttt{invalidateNonce} &
            \texttt{nonce: uint256} &
            Vô hiệu nonce, hủy mọi lệnh dùng nonce đó. \\
        \hline
    \end{tabular}
\end{table}

Hàm \texttt{register} kiểm tra chữ ký của lệnh rồi chuyển tiếp yêu cầu đặt cọc sang \texttt{FillAuction}. Hàm \texttt{executePartialChunk} thực hiện một lần khớp theo trình tự ``kiểm tra - thay đổi trạng thái - tương tác'' (Checks-Effects-Interactions), một quy ước phổ biến để tránh lỗi gọi lại (reentrancy). Cụ thể, trước khi chuyển tài sản, hợp đồng kiểm tra: lệnh chưa bị hủy, chưa bị chuyển sang chế độ dự phòng, nonce chưa bị vô hiệu, filler có đăng ký hợp lệ, và lượng khớp không vượt phần còn lại. Sau khi qua hết các kiểm tra, hợp đồng mới cập nhật phần còn lại rồi mới chuyển tài sản.

Mỗi lần khớp được tính giá độc lập tại block đó theo công thức: số tiền ra bằng lượng khớp nhân với giá tại block đó. Cách này tránh được lỗi mà thiết kế ban đầu gặp phải: tích lũy kết quả toàn bộ lệnh có thể gây tràn số ở lần khớp cuối, khiến lệnh bị kẹt. Ngoài ra, mỗi lần khớp phải đạt mức ra tối thiểu theo tỉ lệ; khi lệnh khớp xong, tổng số tiền ra phải đạt \texttt{minOutputAmount} đã ký.

Giá hiện tại của lệnh tại block $b$ được tính tuyến tính từ giá khởi điểm $P_0$ (\texttt{startPrice}) và mức giảm mỗi block $d$ (\texttt{decayPerBlock}), sàn ở 0:
\begin{equation}
    P(b) = \max\bigl(0,\; P_0 - (b - b_0)\cdot d\bigr)
    \label{eq:decay}
\end{equation}
trong đó $b_0$ là block tại lần khởi tạo gần nhất. Tính an toàn của giá dựa trên hai yếu tố: một là biến thời gian $b$ là \texttt{block.number} do chuỗi cung cấp, một đại lượng tăng đều mà mọi nút đều xác minh được, nên không bên nào có thể chèn giá trị giả; hai là ba tham số $P_0$, $d$, $b_0$ đều nằm trong phần dữ liệu được ký và được cosigner xác thực, nên filler không thể tự đặt giá khởi điểm hay tốc độ giảm giá có lợi cho bản thân. Thứ duy nhất filler còn chọn được là thời điểm khớp; khớp muộn nhận được giá thấp hơn nhưng \texttt{minOutputAmount} đặt mức sàn từ dưới, do đó khoản giá trị mà filler khớp muộn có thể trích xuất bị giới hạn và không thể khiến người dùng nhận dưới mức đã ký.

\begin{lstlisting}[caption={Tính giá hiện tại theo block (DecayCursorLib).}]
function getCurrentPrice(DecayCursor storage c) internal view returns (uint128) {
    uint256 blocksPassed = block.number - c.lastResetBlock;
    uint256 decayed      = blocksPassed * c.decayPerBlock;
    if (decayed >= c.currentStartPrice) return 0;   // cham san
    return c.currentStartPrice - uint128(decayed);
}
\end{lstlisting}

Một lệnh trong hợp đồng có thể ở các trạng thái như ở Hình~\ref{fig:reactor_state}: mới tạo, đang khớp dở, đã khớp xong, hoặc bị hủy / chuyển sang dự phòng. Người dùng cũng có thể chủ động hủy lệnh hoặc vô hiệu một nonce để dừng các lần khớp tiếp theo.

\begin{figure}[H]
    \centering
    \includegraphics[width=0.8\textwidth]{Hinhve/reactor_state.png}
    \caption{Các trạng thái của một lệnh trong hợp đồng PartialFillReactor.}
    \label{fig:reactor_state}
\end{figure}

\section{Hợp đồng FillAuction}
\label{subsec:auction}

\texttt{FillAuction} đảm nhận phần đặt cọc và phạt. Bảng~\ref{table:fa_methods} liệt kê các hàm chính.

\begin{table}[H]
    \centering
    \caption{Các hàm chính của \texttt{FillAuction}.}
    \label{table:fa_methods}
    \begin{tabular}{|p{4cm}|p{4.5cm}|p{5.5cm}|}
        \hline
        \textbf{Hàm} & \textbf{Tham số / Kiểu trả về} & \textbf{Mục đích} \\
        \hline
        \texttt{register} &
            \texttt{orderHash, filler,} \newline \texttt{fillBps, orderSizeEth,} \newline \texttt{deadline} &
            Ghi đăng ký, tính và khóa cọc theo cỡ lệnh và thời gian còn lại. \\
        \hline
        \texttt{onFillSuccess} &
            \texttt{orderHash, filler,} \newline \texttt{filledBps: uint16} &
            Hoàn cọc theo tỉ lệ lượng thực giao / lượng đã cam kết. \\
        \hline
        \texttt{slash} &
            \texttt{orderHash, filler} &
            Phạt filler bỏ dở sau hạn: 90\% vào quỹ, 10\% cho người gọi. \\
        \hline
        \texttt{release}\newline\texttt{Registration} &
            \texttt{orderHash, filler} &
            Hoàn 100\% cọc khi filler không có lỗi (lệnh đã hủy, nonce vô hiệu...). \\
        \hline
        \texttt{withdraw} &
            \texttt{to: address payable} &
            Filler rút ETH đã được ghi nợ (pull payment). \\
        \hline
    \end{tabular}
\end{table}

Khi một filler đăng ký, hợp đồng tính ra số cọc cần nộp và ghi lại đăng ký đó. Số cọc được tính bằng cách quy giá trị phần lệnh về ETH, rồi nhân với một tỉ lệ theo cỡ lệnh và một hệ số theo thời gian còn lại tới hạn (đăng ký càng sát hạn thì hệ số càng cao). Giá trị quy về ETH được lấy từ giá trung bình theo thời gian (TWAP) của pool Uniswap tương ứng.

Cụ thể, số cọc bằng giá trị phần lệnh quy về ETH ($V_{\text{eth}}$) nhân với một tỉ lệ theo cỡ lệnh và một hệ số theo thời gian còn lại:
\begin{equation}
    \text{collateral} = V_{\text{eth}} \cdot \frac{r_s}{10^4} \cdot \frac{m_t}{10^4}
    \label{eq:collateral}
\end{equation}
Hệ số thời gian $m_t$ tăng dần khi đăng ký càng sát hạn, để chặn việc đăng ký sát hạn với cọc rẻ: $m_t = 1{,}0$ khi còn $>50$ block; $1{,}5$ khi còn $>20$; $3{,}0$ khi còn $>5$; và $5{,}0$ khi đã rất sát hạn. Vì cọc tỉ lệ \emph{thẳng} với $V_{\text{eth}}$ (không có chiều tỉ lệ khớp), một cam kết phần lớn không bao giờ rẻ hơn phần nhỏ --- loại bỏ khả năng khai gian cỡ lệnh.

\begin{lstlisting}[caption={Tính cọc động theo cỡ lệnh và thời gian (DynamicStakeLib).}]
function computeCollateral(uint256 notionalEth, uint256 deadline, uint32[4] storage rate)
    internal view returns (uint256)
{
    uint8  sBucket = getOrderSizeBucketETH(notionalEth);  // <1 / <10 / <100 / >=100 ETH
    uint8  tBucket = getTimeBucket(deadline);             // theo so block con lai
    uint256 base   = notionalEth * rate[sBucket] / 10000; // x ti le co lenh
    return base * _getTimeMultiplier(tBucket) / 10000;    // x he so thoi gian
}
\end{lstlisting}

Khi filler khớp xong, hàm \texttt{onFillSuccess} hoàn cọc theo bảng hoàn cọc ở Bảng~\ref{table:refund}. Mức hoàn được tính theo tỉ lệ giữa lượng filler đã giao và lượng đã \emph{cam kết} của chính filler đó, không phải so với cả lệnh. Nhờ vậy, một filler giao đúng phần đã nhận (dù lớn hay nhỏ) được hoàn lại toàn bộ cọc; khi giao thiếu so với cam kết, phần thiếu bị giữ lại làm khoản phạt.

\begin{table}[H]
    \centering
    \caption{Bảng hoàn cọc mặc định (đơn vị: phần vạn của số cọc), theo cỡ lệnh và tỉ lệ giao so với cam kết.}
    \label{table:refund}
    \begin{tabular}{|l|c|c|c|c|c|}
        \hline
        \textbf{Cỡ lệnh \textbackslash{} Tỉ lệ giao} & \textbf{0--2\%} & \textbf{2--10\%} & \textbf{10--30\%} & \textbf{30--70\%} & \textbf{$\geq$70\%} \\
        \hline
        Nhỏ ($<1$ ETH)     & 500 & 1000 & 2500 & 5000 & 10000 \\
        \hline
        Vừa ($1$--$10$)    & 333 & 1000 & 2000 & 5000 & 10000 \\
        \hline
        Lớn ($10$--$100$)  & 100 & 333  & 1000 & 3333 & 10000 \\
        \hline
        Rất lớn ($>100$)   & 100 & 200  & 667  & 2000 & 10000 \\
        \hline
    \end{tabular}
\end{table}

Phần xử lý sau khi hết hạn được thiết kế để công bằng với filler. Sau thời hạn cộng thêm một khoảng chờ (50 block), nếu một đăng ký vẫn chưa tất toán thì có hai khả năng. Nếu lệnh đã được khớp xong, đã bị hủy, hoặc nonce đã bị vô hiệu, thì filler không có lỗi và được hoàn lại toàn bộ cọc qua hàm \texttt{releaseRegistration}. Ngược lại, nếu lệnh vẫn còn phần khớp được mà filler bỏ dở, thì filler bị phạt qua hàm \texttt{slash}: phần lớn cọc (90\%) vào quỹ chung, phần còn lại (10\%) cho người gọi hàm. Cách phân biệt này đã được minh họa ở Hình~\ref{fig:slash_decision} của Chương~\ref{chapter:Related_works}. Toàn bộ việc rút tiền dùng cơ chế ``kéo'' (pull payment): tiền được ghi nợ trước, người nhận tự gọi rút sau, để tránh lỗi khi chuyển tiền.

Bảng hoàn cọc được ghi nhận tại thời điểm đăng ký, nên mọi sửa đổi sau đó không ảnh hưởng tới các filler đã đăng ký trước; số cọc tỉ lệ tuyến tính với khối lượng phần cam kết, loại bỏ khả năng cọc cho phần lớn thấp hơn cọc cho phần nhỏ.

\section{Hợp đồng FallbackExecutor}
\label{subsec:fallback}

\texttt{FallbackExecutor} đảm nhận cơ chế dự phòng. Bảng~\ref{table:fe_methods} liệt kê hàm duy nhất của hợp đồng.

\begin{table}[H]
    \centering
    \caption{Hàm chính của \texttt{FallbackExecutor}.}
    \label{table:fe_methods}
    \begin{tabular}{|p{4cm}|p{4.5cm}|p{5.5cm}|}
        \hline
        \textbf{Hàm} & \textbf{Tham số / Kiểu trả về} & \textbf{Mục đích} \\
        \hline
        \texttt{executeFallback} &
            \texttt{order, sig, cosigSig,} \newline \texttt{router: address,} \newline \texttt{routerData: bytes,} \newline \texttt{minOut: uint256} &
            Đánh dấu lệnh sang chế độ dự phòng, Permit2 kéo phần dư, gọi router trong danh sách trắng. \\
        \hline
    \end{tabular}
\end{table}

Cửa sổ hoạt động của hợp đồng giới hạn trong 10 block trước hạn của lệnh. Khi chạy, hợp đồng xác thực chữ ký của lệnh, rồi bật cờ dự phòng cho lệnh và đưa \texttt{remaining} về 0 --- cả hai thao tác hoàn thành trước khi gọi sang sàn ngoài. Sau đó, mọi lần khớp thường đều bị chặn bởi một trong hai điều kiện độc lập: cờ dự phòng đã bật, hoặc \texttt{remaining} đã bằng 0. Phần khối lượng có thể bị xử lý hai lần ở đây chính là phần tài sản vào chưa khớp của người dùng; nếu vừa bị sàn dự phòng đem đi đổi vừa bị một filler khớp thường rút thêm thì tài sản đó sẽ bị kéo hai lần. Vì cả hai luồng cùng thao tác trên \texttt{remaining}, và biến này được cập nhật \emph{trước} khi tài sản di chuyển, không còn khe gọi lại (reentrancy) để chen vào --- phần khối lượng chưa khớp chỉ có thể do đúng một luồng xử lý.

Sau đó, hợp đồng dùng Permit2 lấy phần còn lại của người dùng, rồi gọi tới một sàn giao dịch để hoán đổi. Kết quả nhận được phải đạt cả mức tối thiểu mà người dùng đã ký lẫn mức tối thiểu do route ước tính; nếu thiếu thì cả giao dịch bị hủy. Nếu route không dùng hết tài sản vào, phần thừa được trả lại cho người dùng.

Về các sàn được hỗ trợ, hợp đồng giữ một danh sách trắng các địa chỉ sàn được phép gọi tới. Đồ án đã cài đặt bộ chuyển đổi cho \textbf{Uniswap} và \textbf{KyberSwap}: bộ chuyển đổi này nằm ở backend, đảm nhận việc gọi API của từng sàn để lấy đường đi tốt nhất rồi mã hóa dữ liệu giao dịch đúng định dạng của sàn đó. Việc bổ sung một sàn mới đòi hỏi viết thêm một bộ chuyển đổi tương tự và thêm địa chỉ sàn vào danh sách trắng, không cần sửa lại hợp đồng.

\section{Các hợp đồng cho giao dịch xuyên chuỗi}
\label{subsec:cc-contract}

Phần xuyên chuỗi dùng một mô hình khác với hợp đồng khóa thời gian dùng chung thường thấy: mỗi phần (slot) của lệnh có một hợp đồng ký quỹ riêng trên cả hai chuỗi, được tạo ra bằng kỹ thuật CREATE2. Trên chuỗi nguồn (ví dụ chuỗi giữ WETH), \texttt{EscrowSrcFactory} đảm nhận việc tạo các hợp đồng ký quỹ \texttt{EscrowSrc}; trên chuỗi đích (ví dụ chuỗi giữ USDC) là \texttt{EscrowDstFactory} tạo ra các \texttt{EscrowDst}. Bảng~\ref{table:cc_methods} liệt kê các hàm chính của cả bốn hợp đồng.

\begin{table}[H]
    \centering
    \caption{Các hàm chính của nhóm hợp đồng xuyên chuỗi.}
    \label{table:cc_methods}
    \begin{tabular}{|p{3.4cm}|p{4.9cm}|p{5.2cm}|}
        \hline
        \textbf{Hàm} & \textbf{Tham số / Kiểu trả về} & \textbf{Mục đích} \\
        \hline
        \multicolumn{3}{|l|}{\textbf{EscrowSrcFactory}} \\
        \hline
        \texttt{fillSlot} &
            \texttt{info, swapperSig,} \newline \texttt{cosignerSig,} \newline \texttt{slotIndex:~uint8,} \newline \texttt{hashlock:~bytes32,} \newline \texttt{proof:~bytes32[]} \newline kèm \texttt{msg.value} (cọc) \newline $\rightarrow$ \texttt{address} &
            Kiểm 2 chữ ký (lần đầu mỗi lệnh) và Merkle proof của slot, tạo EscrowSrc bằng CREATE2, Permit2 kéo tài sản slot vào. \\
        \hline
        \texttt{reopenSlot} &
            \texttt{orderHash, slotIndex} &
            Mở lại slot cho filler mới sau khi filler cũ bỏ. \\
        \hline
        \multicolumn{3}{|l|}{\textbf{EscrowSrc} (mỗi slot một instance)} \\
        \hline
        \texttt{withdraw} &
            \texttt{secret: bytes32} &
            Filler rút tài sản nguồn khi tiết lộ đúng secret. \\
        \hline
        \texttt{cancel} &
            --- &
            Sau T1: hoàn tài sản cho người dùng, tịch thu cọc an toàn của filler. \\
        \hline
        \multicolumn{3}{|l|}{\textbf{EscrowDstFactory}} \\
        \hline
        \texttt{deploy} &
            \texttt{hashlock:~bytes32,} \newline \texttt{filler, token,} \newline \texttt{amount, deadline} \newline $\rightarrow$ \texttt{address} &
            Filler tạo EscrowDst bằng CREATE2, khóa tài sản đích vào hợp đồng mới. \\
        \hline
        \multicolumn{3}{|l|}{\textbf{EscrowDst} (mỗi slot một instance)} \\
        \hline
        \texttt{claim} &
            \texttt{secret: bytes32} &
            Backend/watcher giải phóng tài sản đích khi có secret. \\
        \hline
        \texttt{refund} &
            --- &
            Sau T2: hoàn tài sản về cho filler nếu chưa bị claim. \\
        \hline
    \end{tabular}
\end{table}

Khi một filler nhận một slot, hợp đồng nhà máy kiểm tra hai chữ ký (của người dùng và của cosigner) cùng với bằng chứng Merkle cho slot đó, rồi mới tạo hợp đồng ký quỹ và dùng Permit2 chuyển đúng phần tài sản của slot vào hợp đồng vừa tạo. Filler phải nộp khoản đặt cọc an toàn, khiến hành vi chiếm giữ slot mà không thực hiện trở nên tốn kém. Cây Merkle (mục~\ref{section:2.6}) cho phép xác thực từng slot độc lập với chỉ một chữ ký cho toàn lệnh.

Cụ thể, từ khóa chính $M$ được tính ngoài chuỗi từ khóa gốc và dữ liệu lệnh, mỗi slot $i$ có khóa ($S_i$), hashlock ($h_i$) và lá Merkle riêng:
\begin{align}
    S_i           &= H(M \,\|\, i)                       & \text{(khóa của slot $i$)} \notag\\
    h_i           &= H(S_i)                              & \text{(hashlock)} \notag\\
    \text{leaf}_i &= H\bigl(H(h_i \,\|\, i)\bigr)        & \text{(lá Merkle, băm hai lần)}
    \label{eq:slot-secret}
\end{align}
với $H$ là \texttt{keccak256}. Lá được băm \emph{hai lần} theo chuẩn của thư viện OpenZeppelin~\cite{openzeppelin_contracts} để chống tấn công tiền ảnh thứ hai (xem Phụ lục~\ref{apdx:merkle}). Chỉ gốc Merkle (\texttt{merkleRoot}) được người dùng và cosigner ký rồi đưa lên chuỗi; các $S_i$ không được lưu trữ mà được tính lại khi cần. Khi một filler nhận slot, hợp đồng xác minh cặp $(h_i, i)$ là một lá hợp lệ của cây \emph{trước khi} tiếp cận tài sản, ngăn filler tự chọn hashlock mà bản thân biết trước secret tương ứng.

\begin{lstlisting}[caption={Xác minh lá Merkle của slot trước khi kéo tài sản (EscrowSrcFactory).}]
bytes32 leaf = keccak256(bytes.concat(keccak256(abi.encode(hashlock, slotIndex))));
require(MerkleProof.verify(merkleProof, order.merkleRoot, leaf), "invalid merkle proof");
// ... chi sau dong nay moi Permit2 keo tai san cua slot vao escrow
\end{lstlisting}

Bảng~\ref{table:cc_compare} so sánh mô hình này với cách dùng một hợp đồng khóa thời gian dùng chung.

\begin{table}[H]
    \centering
    \caption{So sánh hai cách làm xuyên chuỗi.}
    \label{table:cc_compare}
    \begin{tabular}{|>{\raggedright\arraybackslash}p{4.2cm}|>{\raggedright\arraybackslash}p{4cm}|>{\raggedright\arraybackslash}p{4cm}|}
        \hline
        \textbf{Tiêu chí} & \textbf{Hợp đồng dùng chung} & \textbf{Hợp đồng riêng mỗi slot} \\
        \hline
        Cô lập sự cố & Một lỗi ảnh hưởng mọi filler & Một hợp đồng hỏng không ảnh hưởng slot khác \\
        \hline
        Tài sản của slot chưa khớp & Nằm trong hợp đồng từ đầu & Vẫn nằm trong ví người dùng \\
        \hline
        Đặt cọc an toàn của filler & Không có & Có, mỗi slot một khoản \\
        \hline
        Dễ kiểm tra trạng thái & Khó & Mỗi hợp đồng tự báo trạng thái riêng \\
        \hline
    \end{tabular}
\end{table}

Một điều kiện thời gian luôn được giữ để bảo vệ filler: hạn trên chuỗi đích (T2) luôn sớm hơn hạn trên chuỗi nguồn (T1). Nhờ vậy, một filler đã bỏ tiền ra trên chuỗi đích nhưng vì lý do nào đó không hoàn tất được thì vẫn lấy lại được tiền của mình trước khi phần tài sản trên chuỗi nguồn bị hoàn cho người dùng. Hình~\ref{fig:sequence_crosschain} và Hình~\ref{fig:crosschain_timing} minh họa luồng 6 bước và ràng buộc thời gian T1/T2.

\begin{figure}[H]
    \centering
    \noindent\makebox[\linewidth][c]{\includegraphics[width=1.3\linewidth]{Hinhve/sequence_crosschain.png}}
    \caption{Sơ đồ trình tự luồng giao dịch xuyên chuỗi (6 bước).}
    \label{fig:sequence_crosschain}
\end{figure}

\begin{figure}[H]
    \centering
    \noindent\makebox[\linewidth][c]{\includegraphics[width=1.3\linewidth]{Hinhve/crosschain_timing.png}}
    \caption{Ràng buộc thời gian T2 < T1 bảo vệ filler khỏi mất cả hai đầu.}
    \label{fig:crosschain_timing}
\end{figure}

\section{Backend và Frontend}
\label{sec:offchain}

Ngoài tầng hợp đồng, hệ thống còn có backend và frontend để kết nối người dùng với các hợp đồng. \textbf{Backend} viết bằng Node.js và TypeScript, dùng cơ sở dữ liệu PostgreSQL. Backend cung cấp các API để tạo và đọc lệnh, giữ vai trò \emph{cosigner} (đồng ký để xác nhận đường giảm giá và gốc cây Merkle), và chạy các thành phần theo dõi (watcher) để tự động kích hoạt cơ chế dự phòng cũng như hoàn tất giao dịch xuyên chuỗi. Backend không giữ tài sản của người dùng: tài sản luôn nằm trong các hợp đồng trên chuỗi, và khi backend gặp sự cố, tiền của người dùng vẫn an toàn. \textbf{Frontend} là giao diện web để người dùng tạo và ký lệnh, theo dõi tiến độ khớp, và xem đường giảm giá; Hình~\ref{fig:frontend} minh họa một số màn hình chính.

\begin{figure}[H]
    \centering
    \begin{subfigure}{0.48\textwidth}
        \centering
        \includegraphics[width=\textwidth]{Hinhve/frontend_1.png}
        \caption{Màn hình tạo lệnh.}
        \label{fig:frontend_create}
    \end{subfigure}
    \hfill
    \begin{subfigure}{0.48\textwidth}
        \centering
        \includegraphics[width=\textwidth]{Hinhve/frontend_2.png}
        \caption{Màn hình theo dõi lệnh và đồ thị đấu giá.}
        \label{fig:frontend_track}
    \end{subfigure}
    \caption{Một số màn hình của giao diện người dùng.}
    \label{fig:frontend}
\end{figure}

\begin{figure}[H]
    \centering
    \begin{subfigure}{0.48\textwidth}
        \centering
        \includegraphics[width=\textwidth]{Hinhve/crosschain_1.png}
        \caption{Màn hình tạo lệnh xuyên chuỗi.}
        \label{fig:crosschain_1}
    \end{subfigure}
    \hfill
    \begin{subfigure}{0.48\textwidth}
        \centering
        \includegraphics[width=\textwidth]{Hinhve/crosschain_2.png}
        \caption{Màn hình theo dõi trạng thái từng slot.}
        \label{fig:crosschain_2}
    \end{subfigure}
    \caption{Giao diện người dùng cho giao dịch xuyên chuỗi.}
    \label{fig:crosschain_ui}
\end{figure}

Việc các hợp đồng và cơ chế vừa trình bày có hoạt động đúng như thiết kế hay không sẽ được kiểm chứng qua các thử nghiệm ở chương tiếp theo.

\end{document}
